% from ega0-17-fr.tex
\oldpage[0\textsubscript{IV-1}]{49}
de b) et d) implique trivialement c), et en vertu de (i), c) entraîne d)~; on
a donc prouvé l'équivalence de b), c), d) et e). En outre, il est clair que b)
entraîne a), les classes des $x_i$ dans
$\mathfrak{m}/\mathfrak{m}^2$ formant alors une base de ce $k$-espace
vectoriel et les classes des $y_i$ dans $\mathfrak{n}/\mathfrak{n}^2$ un
système libre de ce $k'$-espace vectoriel. Reste donc à prouver que a)
entraîne e). Posons $V=\mathfrak{m}/\mathfrak{m}^2$,
$W=\mathfrak{n}/\mathfrak{n}^2$~; il est immédiat que l'on a
$\psi(V\otimes_k k')=(\mathfrak{n}^2+\mathfrak{m}B)/\mathfrak{n}^2$. On a
d'une part, en vertu de (16.3.9)
\begin{equation}
\label{0.17.3.3.3-fr}
\tag{17.3.3.3}
  \dim(B)\leq\dim(A)+\dim(B\otimes_A k)~;
\end{equation}
en second lieu, d'après (16.2.6), on a
\begin{equation}
\label{0.17.3.3.4-fr}
\tag{17.3.3.4}
  \dim(A)\leq\operatorname{rg}_{k'}(V\otimes_k k').
\end{equation}
Enfin, dans l'anneau local $B\otimes_A k=B/\mathfrak{m}B$,
$\mathfrak{n}'=\mathfrak{n}/\mathfrak{m}B$ est l'idéal maximal, $k'$ le corps
résiduel et $\mathfrak{n}'/\mathfrak{n}'^2$ est isomorphe à
$\mathfrak{n}/(\mathfrak{n}^2+\mathfrak{m}B)
=W/\psi(V\otimes_k k')$~; on a donc, d'après (16.2.6)
\begin{equation}
\label{0.17.3.3.5-fr}
\tag{17.3.3.5}
  \dim(B\otimes_A k)
  \leq\operatorname{rg}_{k'}W
  -\operatorname{rg}_{k'}\psi(V\otimes_k k').
\end{equation}
Enfin, puisque $B$ est supposé régulier, on a
$\dim(B)=\operatorname{rg}_{k'}W$ (17.1.1)~; on conclut donc de (17.3.3.3),
(17.3.3.4) et (17.3.3.5) que l'on a
\[
  \operatorname{rg}_{k'}W
  \leq\dim(A)+\dim(B\otimes_A k)
  \leq\operatorname{rg}_{k'}W
  +\operatorname{rg}_{k'}(V\otimes_k k')
  -\operatorname{rg}_{k'}\psi(V\otimes_k k')
\]
et dire que les deux termes extrêmes de cette inégalité sont égaux signifie
que $\psi$ est \emph{injectif}. La condition a) entraîne donc nécessairement
que dans chacune des relations (17.3.3.3), (17.3.3.4) et (17.3.3.5), les deux
membres sont \emph{égaux}~; or, l'égalité dans (17.3.3.4) (resp. (17.3.3.5))
signifie que $A$ (resp. $B\otimes_A k$) est \emph{régulier} (17.1.1)~; on a
donc bien prouvé que a) entraîne e).
\end{proof}

\begin{proposition}[17.3.4]
\label{0.17.3.4-fr}
Soient $A$ un anneau local régulier de dimension $n$, $\mathfrak{m}$ son idéal
maximal. Pour tout $A$-module non nul de type fini $M$, on a
\begin{equation}
\label{0.17.3.4.1-fr}
\tag{17.3.4.1}
  \operatorname{prof}(M)+\operatorname{dim.proj}(M)=n.
\end{equation}
\end{proposition}

\begin{proof}
Raisonnons par récurrence sur $m=\operatorname{prof}(M)$. Si $m=0$, on sait
(16.4.6, (i)) qu'il existe un sous-$A$-module $N$ de $M$ isomorphe à
$k=A/\mathfrak{m}$~; appliquant la suite exacte des Tor à la suite exacte
$0\longrightarrow N\longrightarrow M\longrightarrow M/N\longrightarrow0$,
on obtient une suite exacte
\[
  \operatorname{Tor}^A_{n+1}(M/N,k)\longrightarrow
  \operatorname{Tor}^A_n(N,k)\longrightarrow
  \operatorname{Tor}^A_n(M,k),
\]
et l'hypothèse que $A$ est régulier entraîne
$\operatorname{Tor}^A_{n+1}(M/N,k)=0$ ((17.3.1.1) et (17.2.6)), donc
$\operatorname{Tor}^A_n(k,k)$ est isomorphe à un sous-$A$-module de
$\operatorname{Tor}^A_n(M,k)$~; appliquant de nouveau (17.3.1.1) et (17.2.6),
on voit que $\operatorname{Tor}^A_n(M,k)\neq0$, donc
$\operatorname{dim.proj}(M)\geq n$ (17.2.6)~; mais comme
$\operatorname{dim.proj}(M)\leq n$ par (17.3.1.1), on a bien
$\operatorname{dim.proj}(M)=n$.

Supposons maintenant $m>0$, et soit $x$ un élément $M$-régulier appartenant à
$\mathfrak{m}$~; on sait alors (16.4.6, (i)) que l'on a
$\operatorname{prof}(M/xM)=m-1$, et d'autre part (17.2.7)
$\operatorname{dim.proj}(M/xM)=\operatorname{dim.proj}(M)+1$~; l'hypothèse de
récurrence prouve aussitôt la relation (17.3.4.1).
\end{proof}

